The \peripheral{I2Cx} peripherals are general-purpose I\textsuperscript{2}C interfaces, each able to act as a bus master or as an addressed slave. Master and slave are driven through two separate sets of register bits sharing one pin pair, so a port can be configured for either role (or, with both enabled, participate in a multi-master bus). The main features are listed below:

\begin{itemize}
	\item Master transmitter and master receiver modes, with START, repeated-START and STOP generation (\bitfield{I2CMST}, \bitfield{I2CMSP}) and a programmable bus clock divider (\bitfield{I2CMDIV})
	\item Multi-master arbitration: losing the bus releases it and raises \bitfield{I2CMARB}
	\item Slave transmitter and slave receiver modes with a 7-bit address (\register{I2CxAR}), a per-bit address mask (\register{I2CxAMR}) for responding to a range of addresses, and an optional general-call response (\bitfield{I2CGCE})
	\item Optional slave clock stretching (\bitfield{I2CSCS}): the port holds \pin{SCLx} low through the ACK bit until firmware releases it, trading bus latency for lossless flow control
	\item Separate transmit and receive data registers for each role (\register{I2CxMTX}/\register{I2CxMRX}, \register{I2CxSTX}/\register{I2CxSRX})
	\item Individually maskable interrupt sources for transfer complete, arbitration loss, NACK received, address match, START and STOP detection, and receive overrun
\end{itemize}

\subsection{The Wire Protocol}

Both roles speak the same two-wire protocol. \pin{SDAx} and \pin{SCLx} are open-drain and require external pull-up resistors: a device drives a bit low by pulling the line down and releases it to let the pull-up present a one. Because every device can only pull down, two devices driving at once cannot damage each other --- that is what makes both multi-master arbitration and slave clock stretching possible on the same pair of wires.

A transfer is framed by two conditions that are deliberately illegal during data: \textbf{START} is a falling edge on \pin{SDAx} while \pin{SCLx} is high, and \textbf{STOP} is a rising edge on \pin{SDAx} while \pin{SCLx} is high. Between them, \pin{SDAx} may only change while \pin{SCLx} is low, so every data bit is stable across the whole \pin{SCLx} high period and any receiver may sample it there. After the address byte and after each data byte, the transmitter releases \pin{SDAx} for one clock and the receiver answers by pulling it low for an ACK or leaving it high for a NACK. Figure \ref{fig:i2c-transaction} shows one complete byte write.

% Generated by LatexUserGuide.GenerateI2cTransactionDiagram (make chip).
\begin{figure}[htpb]
	\centering
	\input{include/I2cTransactionDiagram.tex}
	\caption{One I2C byte write on the wire: START, the 7-bit address and direction bit, the addressed device's ACK, one data byte, its ACK, then STOP. START and STOP are the two transitions of \pin{SDAx} that occur while \pin{SCLx} is high; every data bit is stable across the \pin{SCLx} high period.}
	\label{fig:i2c-transaction}
\end{figure}

The address byte carries the 7-bit slave address in its most significant bits and the direction in its least significant bit: 0 selects a write (master transmitter), 1 a read (master receiver). A master that wants to change direction or address without giving up the bus issues a repeated START instead of a STOP.

\subsection{Clock Domain}

The \peripheral{I2Cx} cores are clocked from SMCLK, and \bitfield{I2CMDIV} divides that down to the bus rate. The boot ROM may leave SMCLK sourced from LFXT (32.768\,kHz), which would put the bus in the low-kilohertz range --- configure the SMCLK source through \register{SYSCLKCR} (see the \peripheral{SYSTEM} chapter) before relying on any I\textsuperscript{2}C timing, rather than inheriting whatever the boot path left behind.

The step-by-step register sequences for each of the four roles are given with the register descriptions below.
